% GENERATED FILE. Edit canonical lessons, then run npm run generate:books.
% canonical-source-hash: fnv1a64:3a01e0c6491b544c
% canonical-lessons: JA-C352-koboreru, JA-C352-hodokeru, JA-C352-hikkurikaeru, JA-C352-zureru, JA-C352-dougu

\chapter{Little Accidents}
\label{ch:ja-little-accidents}
\hlchaptermodality{352}

\begin{chapteropening}
\textbf{By the end of this chapter:} \emph{I can say what to do if something spills, tears or falls over.}

Complete the last lesson of chapter 352: I can say what to do if something spills, tears or falls over.
\end{chapteropening}

\section[\texorpdfstring{\ja{こぼれる} (koboreru)}{こぼれる (koboreru)}]{\ja{こぼれる} (koboreru) — to be spilled; to overflow}

\label{lesson:JA-C352-koboreru}



\begin{quote}
Before the new one: say the Japanese for whether someone is safe, then the Japanese for to be on the safe side.
\end{quote}

\subsection*{You'll want to know: \ja{こぼれる}}
\textbf{\ja{こぼれる}} — \emph{koboreru} — \textquotedblleft{}to be spilled; to overflow\textquotedblright{}.

\begin{quote}
\textbf{\ja{みずが} \ja{こぼれたら、ゆかが} \ja{ぬれます。}} — \emph{mizu ga koboretara, yuka ga nuremasu} — \textquotedblleft{}If the water spills, the floor gets wet.\textquotedblright{}
\end{quote}

Polite \textbf{\ja{こぼれます}}, \emph{koboremasu}; an ichidan verb. Past (polite): \textbf{\ja{ぎゅうにゅうが} \ja{こぼれました。}}, \emph{gyūnyū ga koboremashita}.

\begin{cousinweb}
\textbf{\ja{こぼれる}} (\emph{koboreru}) is a native Japanese word (wago), inherited rather than borrowed from Chinese, and its older origin is not settled. JMdict notes that it is usually written in kana (JMdict).
\end{cousinweb}

\begin{grammarlens}[title={What you've built}]
One more word for saying what happens if.
\end{grammarlens}

\subsection*{Your turn}
\begin{itemize}
\raggedright
  \item \emph{Say it:} \emph{koboreru}
  \item \emph{Say it:} \emph{koboreru}, once more
  \item \emph{Say it:} \emph{mizu ga koboretara, yuka ga nuremasu} — If the water spills, the floor gets wet.
  \item \emph{Recall:} say the Japanese for whether someone is safe, then the Japanese for to be on the safe side, then say \emph{koboreru} again
\end{itemize}

\begin{tcolorbox}[breakable,colback=teal!4,colframe=teal!35!black,title={Before you move on}]
What does \ja{こぼれる} mean? (\textquotedblleft{}To be spilled; to overflow\textquotedblright{}.) Say it once more.
\end{tcolorbox}
\section[\texorpdfstring{\ja{ほどける} (hodokeru)}{ほどける (hodokeru)}]{\ja{ほどける} (hodokeru) — to come undone; to come untied}

\label{lesson:JA-C352-hodokeru}



\begin{quote}
Before the new one: say the Japanese for to be on the safe side, then the Japanese for to be spilled.
\end{quote}

\subsection*{You'll want to know: \ja{ほどける}}
\textbf{\ja{ほどける}} — \emph{hodokeru} — \textquotedblleft{}to come undone; to come untied\textquotedblright{}.

\begin{quote}
\textbf{\ja{くつの} \ja{ひもが} \ja{ほどけたら、むすんで} \ja{ください。}} — \emph{kutsu no himo ga hodoketara, musunde kudasai} — \textquotedblleft{}If your shoelace comes undone, please tie it.\textquotedblright{}
\end{quote}

Polite \textbf{\ja{ほどけます}}, \emph{hodokemasu}; an ichidan verb. Past (polite): \textbf{\ja{ひもが} \ja{ほどけました。}}, \emph{himo ga hodokemashita}.

\begin{cousinweb}
\textbf{\ja{ほどける}} (\emph{hodokeru}) is a native Japanese word (wago), inherited rather than borrowed from Chinese, and its older origin is not settled. JMdict notes that it is usually written in kana (JMdict).
\end{cousinweb}

\begin{grammarlens}[title={What you've built}]
One more word for saying what happens if.
\end{grammarlens}

\subsection*{Your turn}
\begin{itemize}
\raggedright
  \item \emph{Say it:} \emph{hodokeru}
  \item \emph{Say it:} \emph{hodokeru}, once more
  \item \emph{Say it:} \emph{kutsu no himo ga hodoketara, musunde kudasai} — If your shoelace comes undone, please tie it.
  \item \emph{Recall:} say the Japanese for to be on the safe side, then the Japanese for to be spilled, then say \emph{hodokeru} again
\end{itemize}

\begin{tcolorbox}[breakable,colback=teal!4,colframe=teal!35!black,title={Before you move on}]
What does \ja{ほどける} mean? (\textquotedblleft{}To come undone; to come untied\textquotedblright{}.) Say it once more.
\end{tcolorbox}
\section[\texorpdfstring{\ja{ひっくりかえる} (hikkurikaeru)}{ひっくりかえる (hikkurikaeru)}]{\ja{ひっくりかえる} (hikkurikaeru) — to tip over; to turn over}

\label{lesson:JA-C352-hikkurikaeru}



\begin{quote}
Before the new one: say the Japanese for to be spilled, then the Japanese for to come undone.
\end{quote}

\subsection*{You'll want to know: \ja{ひっくりかえる}}
\textbf{\ja{ひっくりかえる}} — \emph{hikkurikaeru} — \textquotedblleft{}to tip over; to turn over\textquotedblright{}.

\begin{quote}
\textbf{\ja{かびんが} \ja{ひっくりかえったら、みずが} \ja{こぼれます。}} — \emph{kabin ga hikkurikaettara, mizu ga koboremasu} — \textquotedblleft{}If the vase tips over, the water spills.\textquotedblright{}
\end{quote}

Polite \textbf{\ja{ひっくりかえります}}, \emph{hikkurikaerimasu}; a godan verb. Past (polite): \textbf{\ja{かびんが} \ja{ひっくりかえりました。}}, \emph{kabin ga hikkurikaerimashita}.

\begin{cousinweb}
\textbf{\ja{ひっくりかえる}} (\emph{hikkurikaeru}) is a native Japanese word (wago), inherited rather than borrowed from Chinese: it is built on \emph{kaeru}, \textquotedblleft{}to turn over\textquotedblright{}, the intransitive partner of \emph{kaesu}, \textquotedblleft{}to give back\textquotedblright{}, with a native first part, \emph{hikkuri}, and its kanji was attached later as a kun'yomi (JMdict).
\end{cousinweb}

\begin{grammarlens}[title={What you've built}]
One more word for saying what happens if.
\end{grammarlens}

\subsection*{Your turn}
\begin{itemize}
\raggedright
  \item \emph{Say it:} \emph{hikkurikaeru}
  \item \emph{Say it:} \emph{hikkurikaeru}, once more
  \item \emph{Say it:} \emph{kabin ga hikkurikaettara, mizu ga koboremasu} — If the vase tips over, the water spills.
  \item \emph{Recall:} say the Japanese for to be spilled, then the Japanese for to come undone, then say \emph{hikkurikaeru} again
\end{itemize}

\begin{tcolorbox}[breakable,colback=teal!4,colframe=teal!35!black,title={Before you move on}]
What does \ja{ひっくりかえる} mean? (\textquotedblleft{}To tip over; to turn over\textquotedblright{}.) Say it once more.
\end{tcolorbox}
\section[\texorpdfstring{\ja{ずれる} (zureru)}{ずれる (zureru)}]{\ja{ずれる} (zureru) — to slip out of place}

\label{lesson:JA-C352-zureru}



\begin{quote}
Before the new one: say the Japanese for to come undone, then the Japanese for to tip over.
\end{quote}

\subsection*{You'll want to know: \ja{ずれる}}
\textbf{\ja{ずれる}} — \emph{zureru} — \textquotedblleft{}to slip out of place\textquotedblright{}.

\begin{quote}
\textbf{\ja{めがねが} \ja{ずれたら、なおして} \ja{ください。}} — \emph{megane ga zuretara, naoshite kudasai} — \textquotedblleft{}If your glasses slip, push them back.\textquotedblright{}
\end{quote}

Polite \textbf{\ja{ずれます}}, \emph{zuremasu}; an ichidan verb. Past (polite): \textbf{\ja{めがねが} \ja{ずれました。}}, \emph{megane ga zuremashita}.

\begin{cousinweb}
\textbf{\ja{ずれる}} (\emph{zureru}) is a native Japanese word (wago), inherited rather than borrowed from Chinese, and its older origin is not settled; JMdict gives it no kanji spelling (JMdict).
\end{cousinweb}

\begin{grammarlens}[title={What you've built}]
One more word for saying what happens if.
\end{grammarlens}

\subsection*{Your turn}
\begin{itemize}
\raggedright
  \item \emph{Say it:} \emph{zureru}
  \item \emph{Say it:} \emph{zureru}, once more
  \item \emph{Say it:} \emph{megane ga zuretara, naoshite kudasai} — If your glasses slip, push them back.
  \item \emph{Recall:} say the Japanese for to come undone, then the Japanese for to tip over, then say \emph{zureru} again
\end{itemize}

\begin{tcolorbox}[breakable,colback=teal!4,colframe=teal!35!black,title={Before you move on}]
What does \ja{ずれる} mean? (\textquotedblleft{}To slip out of place\textquotedblright{}.) Say it once more.
\end{tcolorbox}
\section[\texorpdfstring{\ja{どうぐ} (dōgu)}{どうぐ (dōgu)}]{\ja{どうぐ} (dōgu) — a tool; an instrument}

\label{lesson:JA-C352-dougu}



\begin{quote}
Before the new one: say the Japanese for to tip over, then the Japanese for to slip out of place.
\end{quote}

\subsection*{You'll want to know: \ja{どうぐ}}
\textbf{\ja{どうぐ}} — \emph{dōgu} — \textquotedblleft{}a tool; an instrument\textquotedblright{}.

\begin{quote}
\textbf{\ja{どうぐが} \ja{こわれたら、かいしゃに} \ja{れんらくして} \ja{ください。}} — \emph{dōgu ga kowaretara, kaisha ni renraku shite kudasai} — \textquotedblleft{}If a tool breaks, please let the company know.\textquotedblright{}
\end{quote}

\begin{cousinweb}
\textbf{\ja{どうぐ}} (\emph{dōgu}) is Sino-Japanese (kango): two Chinese characters read with their on'yomi, readings borrowed from Chinese. They are \emph{dō}, \textquotedblleft{}road-way\textquotedblright{} (Middle Chinese \emph{dawX}) and \emph{gu}, \textquotedblleft{}tool\textquotedblright{} (\emph{gjuH}) (KANJIDIC2; Middle Chinese in Baxter's transcription).
\end{cousinweb}

\begin{grammarlens}[title={What you've built}]
One more word for saying what happens if.
\end{grammarlens}

\subsection*{Your turn}
\begin{itemize}
\raggedright
  \item \emph{Say it:} \emph{dōgu}
  \item \emph{Say it:} \emph{dōgu}, once more
  \item \emph{Say it:} \emph{dōgu ga kowaretara, kaisha ni renraku shite kudasai} — If a tool breaks, please let the company know.
  \item \emph{Recall:} say the Japanese for to tip over, then the Japanese for to slip out of place, then say \emph{dōgu} again
\end{itemize}

\begin{tcolorbox}[breakable,colback=teal!4,colframe=teal!35!black,title={Before you move on}]
What does \ja{どうぐ} mean? (\textquotedblleft{}A tool; an instrument\textquotedblright{}.) Say it once more.
\end{tcolorbox}
